\chapter{条件视觉生成}
\label{lab:E50_cv_diffusion}

\noindent\textbf{配套文件：}\path{notebook/E50_cv_diffusion.ipynb}\par
\noindent\textbf{教材关联：}\bookxrefshort{ch:diffusion-probability}、\bookxrefshort{ch:conditional-vision}。

\section*{实验目标}
核对前向噪声过程、噪声预测和反向采样的参数化。

\section*{准备及定位}
先阅读本册实验总则，确认Notebook中的依赖与资产单元。原理计算、库接口迁移和真实模型训练可能具有不同资源要求，应分别记录设备、版本及运行范围。

源码定位可从下列函数开始；应结合前置数据与配置单元阅读。
\begin{itemize}
\item \texttt{my\_linear\_beta\_schedule}
\item \texttt{my\_extract}
\item \texttt{my\_add\_noise}
\end{itemize}

\section*{操作及观察}
\begin{enumerate}
\item 固定随机种子，比较不同时间步的加噪结果和噪声计划。
\item 核对时间嵌入、噪声预测器的输入输出及训练目标。
\item 比较有条件与无条件预测，记录CFG系数和VAE潜变量尺度。
\item 检查反向均值与采样器接口，迁移Diffusers时核对预测类型、时间步和调度器。
\end{enumerate}

\section*{结果判读及提交}
提交噪声轨迹、损失定义及采样配置。单次生成观感不构成分布质量评测；不同预测参数化不能直接共用同一反向公式。

提交时附上所执行单元范围、环境与制品身份、原始结果及失败说明。将未运行项目明确列出，不能用Notebook中已有输出替代本次执行证据。

相关方法的原始定义与适用前提见\cite{ho2020ddpm}；本实验的运行结果仍须按实际配置单独验证。
